\documentclass[12pt]{article}
\usepackage[utf8]{inputenc}
\usepackage[T1]{fontenc}
\usepackage{hyperref}
\title{Tarjumo Language Resources for K-12 Curriculum Development}
\author{Compiled by Claude Sonnet 5 (claude-sonnet-5) for Andr\'as Kornai}
\date{}
\begin{document}
\maketitle

\section{Introduction}
Tarjumo (ISO 639-3 txj) is a classical, liturgical language that scholars identify with Old Kanembu. It was used by Muslim scholars in the Kanem--Bornu area and is written in Ajami (Arabic-based) script. It has no known first-language speakers and no ethnic community. The sources I found say it was documented in 2005 among Borno scholars in northeastern Nigeria, who use it for religious recitation and commentary on Arabic texts. It is unintelligible to speakers of modern Kanuri.

\section{Primary Educational Resources}
I found no primers, textbooks or school materials.

\section{Linguistic and Reference Materials}
\begin{itemize}
\item Tarjumo language (Wikipedia): \url{https://en.wikipedia.org/wiki/Tarjumo_language}
\end{itemize}
This is the only source I found. I did not find a grammar or a dictionary.

\section{Religious and Community Resources}
The language lives only in the Islamic scholarly tradition of Borno, in recitation and commentary (tarjuma) on Arabic religious texts. I found no published texts online.

\section{Digital and Audio Resources}
I found none.

\section{Curriculum Development Context}
Little exists. A language-medium K-12 curriculum is not feasible. At most, a heritage unit could introduce the Kanem--Bornu scholarly tradition, with the guidance of Borno scholars. Primary documentation would have to be sought from the researchers named in the Wikipedia article, but I have not verified what they have published.
\end{document}